\begin{cabstract}

在灾后道路排查等无地面网络覆盖的场景中，多架无人机需要从不同视角协同观测现场，并经多跳中继将数据回传至地面中心，供其在规定时限内综合研判目标状态。受机载电量和无线带宽限制，直接回传原始图像开销过大，利用端侧视觉语言模型将画面提炼为紧凑的语义消息成为可行途径。然而，在通信电量、单次任务消息长度和任务时限三方面约束下，端到端回传在网络转发、内容筛选与可靠交付三个环节仍存在突出问题。围绕这些问题，本文开展了三方面研究工作。

（1）在多跳无人机网络中，节点高速移动导致链路频繁变化，而机载通信电量有限。转发时若只看链路速率或地理距离，流量容易集中到少数位置有利的中继上，使其过早耗尽电量而断链；若强行让所有节点平均分担能耗，又会迫使数据走较差的链路，拉低整体吞吐量，难以兼顾传输效率与网络续航。为此，本文提出基于李雅普诺夫优化的分布式路由与功率控制方法 LEADR。该方法将数据队列积压、累计耗电量和虚拟能量队列一并纳入每个时隙的优化目标，把长期电量限制转化为节点本地的在线决策。各节点结合当前信道与干扰直接算出最优发射功率，再综合邻居剩余电量与到中心的地理进度挑选下一跳，同时依据节点间的相对运动预测调整邻居探测周期。在不同无人机数量、通信范围和飞行速度的仿真实验中，LEADR 均未出现节点提前耗尽电量的情况，同时取得了最高的转发吞吐量与最低的丢包率。

（2）链路连通后，单次任务允许发送的消息总长度依然有限，无法把所有观测内容全部回传。多架无人机的视野往往存在重叠，部分视角还可能被遮挡；同时，同一段观测既可以只发简短的类别标签，也可以发附带依据的详细描述，前者占带宽少但信息有限，后者细节丰富但传输耗时长、更容易丢包。如果只按单条消息的置信度来挑选，不仅容易重复发送已有内容、选入受遮挡视角的错误判断，也无法随中心已收信息动态调整。为此，本文提出基于强化学习的语义消息选择方法 EviComRL。该方法先用参数固定的视觉语言模型离线生成各视角不同详略程度的候选消息，并结合传输开销初筛缩小候选范围；在线选择时，通过策略网络综合各节点的观测可靠程度、中心当前判断、已选消息以及链路状态，屏蔽已选或超预算的候选，逐条决定发送哪条消息及何时停止，并采用近端策略优化（PPO）算法进行训练。在五级证据可见性分类的基准实验中，其准确率与宏平均 F1 分别达到 58.17\% 和 50.58\%，在词元开销较低的情况下，分类效果明显好于全量发送和单步排序等对比方法。

（3）选好的消息进入无线信道后，既可能因信道波动发生丢包，地面中心发回的确认消息也可能延迟到达甚至丢失。发送端收不到确认时，无法分辨消息是传输失败还是已经送达，此时立刻重传会白白耗电并挤占后续任务资源，一味等待又容易错过截止时间；此外，即便补发了消息，只要它改变不了地面中心最终的类别判断，这次传输就起不到实际作用。为此，本文提出一致性感知的信念反馈控制方法 C-BFC。该方法先用轻量级神经网络估计不同接收消息组合的判断准确程度，再将预测类别相同的组合估值取均值以保持一致，避免模型因置信度差异冗余发送消息。实际传输时，发送端根据已收到的累计确认和此后的发送记录，推算中心当前收到各条消息的概率，并用虚拟能量队列跟踪前面任务的耗电情况以控制当前支出。在剩余时间内，该方法向前预估两步发送的回报，并与等待确认的预期回报作比较，据此决定发送新消息、重传、等待还是停止。在信道波动、反馈延迟、任务时限和节点电量等 17 组实验条件下，C-BFC 的平均按时正确完成率达到 49.84\%，单位正确任务能耗降为 3.615 mJ，在完成率略占优势的同时，能耗较同环境下重新训练的强化学习方法降低了 14\% 以上。
\end{cabstract}

\begin{eabstract}

In post-disaster road assessment without terrestrial networks, multiple unmanned aerial vehicles (UAVs) observe target areas from different views and relay data across multiple hops to a ground center for timely judgment. Limited onboard energy and bandwidth make raw image backhaul too costly, so distilling local observations into compact semantic messages via onboard vision-language models offers a practical alternative. Under energy, message-length, and deadline constraints, however, end-to-end delivery still faces major difficulties in forwarding, content selection, and reliable delivery. Addressing these issues, this thesis carries out three lines of research.

First, in multi-hop UAV networks, high mobility causes frequent link changes while communication energy is limited. Forwarding solely by link rate or geographic distance concentrates traffic on a few well-positioned relays and drains their batteries early, whereas forcing uniform energy sharing diverts traffic onto poor links and lowers throughput, making it hard to balance efficiency with network lifetime. To address this trade-off, a Lyapunov-based decentralized routing and power control method (LEADR) is proposed. It incorporates queue backlogs, cumulative energy consumption, and virtual energy queues into per-slot optimization, turning long-term energy limits into local online decisions. Each node computes optimal transmit power from channel and interference states, selects next hops by neighbor residual energy and geographic progress toward the center, and adapts neighbor discovery intervals via relative motion prediction. Across simulations with varying UAV counts, communication ranges, and flight speeds, LEADR prevents early battery depletion throughout the mission while achieving the highest forwarding throughput and lowest packet loss.

Second, once links are established, per-task message length budgets remain too tight to send all observations. UAV views often overlap and some may be occluded. Meanwhile, the same observation can be sent as a short class label saving bandwidth with less detail, or as a detailed description with supporting context that takes longer to transmit and risks higher loss. Selecting by single-message confidence alone risks sending redundant content or misleading occluded views and cannot adapt to what the center already has. Accordingly, a reinforcement-learning-based semantic message selection method (EviComRL) is developed. It uses a frozen vision-language model offline to generate candidate messages of different detail levels per view and narrows the pool by transmission cost. Online, a policy network combines node reliability, current center judgment, selected messages, and link states, masks selected or over-budget candidates, and sequentially decides which message to send and when to stop, trained via proximal policy optimization (PPO). In benchmark five-level visibility classification experiments, it reaches 58.17\% accuracy and 50.58\% macro F1, clearly outperforming full transmission and single-step ranking baselines with low token usage.

\begingroup
\clubpenalty=10000
\widowpenalty=10000
Third, once selected messages enter wireless channels, channel fluctuations can drop packets, while ground acknowledgments may be delayed or lost. Without acknowledgments, the sender cannot tell whether a message failed or arrived; immediate retransmission wastes energy and squeezes resources for later tasks, while waiting too long risks missing the deadline. Moreover, resending a message adds no value if it cannot change the center's final predicted class. To handle this, a consistency-aware belief feedback control method (C-BFC) is proposed. It trains a lightweight neural network to estimate the accuracy of received message subsets and averages estimates across subsets predicting the same class, preventing confidence gaps from triggering unnecessary sends that do not alter final decisions. During delivery, the sender infers reception probabilities from cumulative acknowledgments and send logs, while tracking prior energy usage via virtual energy queues to restrain current spending. Within the remaining time, it compares the estimated return of two-step lookahead sending against waiting for acknowledgments to choose new sends, retransmissions, waiting, or stopping. Across 17 experimental settings covering channel fluctuation, feedback delay, task deadlines, and node energy, C-BFC achieves an average on-time correct completion rate of 49.84\% at 3.615 mJ per correct task, cutting energy by over 14\% versus retrained reinforcement-learning methods with slightly higher completion.
\par\endgroup
\end{eabstract}